% Compact evaluation: one validation table and two figures.
% Requires graphicx, amsmath, amssymb, booktabs.
% In the main paper: \newcommand{\quclevalpath}{evaluation}
% Add evaluation/references to the main paper's bibliography database list.
\providecommand{\quclevalpath}{.}
\input{\quclevalpath/numbers.tex}

\section{Experimental Validation and Analysis}
\label{sec:eval}
We ask two questions: how does coupled execution change completion time,
and how much error results from simplifying that execution? We first
validate correctness, then answer these questions with two controlled studies.

\subsection{Setup and Validation}
\label{subsec:setup}
The composed workload performs swapping at R and uses the resulting A--B
pair for teleportation from A to B. Inputs and elementary EPRs are created
at $t=0$. Native NetSquid channels deliver remote halves to A and B in
0.8 and 1.2~ms. R has ideal instantaneous knowledge of this delivery;
after the swap correction, A instead learns pair readiness through an
actual B--R--A UDP packet. Swap and teleportation result packets traverse
R--B and A--R--B. R and A have separate BSM processors; corrections share B.

Classical links operate at 10~Mb/s. Native CNOT/H/M/M instructions take
0.6/0.2/0.4/0.4~ms; X-or-I and Z-or-I correction slots take 0.25~ms each.
Native memory noise uses $T_1/T_2=20/10$~ms in idle and active storage.
Quantum transit is lossless, with zero depolarization and no memory
T1/T2 noise; destination-memory aging begins on placement.
These are controlled parameters, not hardware calibration or a model
of heralded entanglement generation.

\input{\quclevalpath/physical-model.tex}

\input{\quclevalpath/table_validation.tex}

Six Pauli eigenstates, including both $|+i\rangle$ and $|-i\rangle$,
cover all 16 joint swap/teleport measurement outcomes in 96 noiseless
executions. The minimum fidelity is numerically unity. Twelve noisy runs
cover all six inputs at A--R delays of 0.2 and 1~ms. State errors against
an independent two-stage NetSquid reference are shown in
Table~\ref{tab:qucl-validation}. Checks also verify same-object pair
handoff, single consumption, and packet-reception gating of subsequent
operations. State comparisons use observed operation timestamps;
separate per-hop FIFO and processor checks validate causal timing.
The 500-request, zero-error FIFO result in the table is a separate frozen
v4 mixed-protocol regression, not an additional composed-workflow sample.

\subsection{End-to-End Execution}
\label{subsec:coupled}
Figure~\ref{fig:latency} compares four execution conditions at each
A--R propagation delay. R--B propagation remains 0.2~ms; the first start
is at 1~ms. The three delays, two start intervals, two R--B background
loads, and four periodic-traffic phases yield 48 runs and 192 chains.
Inputs are $|+i\rangle$. Latency is final teleportation correction
completion minus local workflow start.

At the 0.8-ms interval without background traffic, mean latency rises
from \quclChainLowLatency to \quclChainMidLatency and
\quclChainHighLatency~ms as A--R delay increases from 0.2 to 1 and 5~ms.
For an isolated workflow with unchanged waits, A--R is crossed by both
readiness and teleportation-result signaling, giving
$\Delta T=2\Delta d_{AR}$. The isolated 0.2-to-1-ms comparison confirms
unchanged swap checkpoints, readiness shifted by 0.8~ms, and completion
shifted by 1.6~ms.

\begin{figure*}[t]
\centering
\includegraphics[width=\textwidth]{\quclevalpath/figures/fig3_end_to_end_latency.pdf}
\caption{Direct comparison of composed-workflow latency. Color denotes
start interval; hatching denotes R$\to$B background load 0.7. Bar heights
and labels are means. Whiskers show observed minima/maxima over four
chain positions and four traffic phases per cell, not confidence
intervals. All conditions share one zero-based axis.}
\label{fig:latency}
\end{figure*}

The delays add, but their components are causally dependent: quantum
completion determines packet generation, and packet arrival determines
eligibility for the next quantum operation. In a separate 24-run
sensitivity at load 0.7, halving quantum RA/RB delivery times reduces
latency by \quclFastLatencyReduction~ms and packet queueing by
\quclFastQueueReduction~ms. Doubling them increases latency by
\quclSlowLatencyIncrease~ms while packet queueing decreases by
\quclSlowQueueReduction~ms. This fixed-ratio sensitivity changes readiness
and remote-memory placement together; it does not isolate a scheduler delay.

Because resources are created at $t=0$, start-interval comparisons change
both contention and storage age. The finite phase sweep is not a population
estimate; at zero background load, phase repetitions are identical conditions.

\subsection{Errors from Execution Simplification}
\label{subsec:abstraction}
The second study uses a separate frozen v4 provisioned swapping workload.
Its correction is 50~$\mu$s, below the 88-$\mu$s result serialization
time, to exclude B contention. Fidelity targets the corrected A--B Bell
pair, rather than the teleported input of the first study.

Full Coupling retains readiness, shared processors, native timed quantum
operations, and actual result packets. No R Contention removes only R's
capacity constraint and generates actual packets at the new completion
times. Fixed Classical Delay replaces result transport with a mean
calibrated from separate traffic phases. Decoupled Timing combines R FIFO
waits with No R Contention packet delays without regenerating packets
at shifted times; it predicts no quantum state.

Three loads, two intervals, 32 test traffic phases, and two quantum seeds
give 384 paired cases and 1,536 transactions per model. The 95\% bootstrap
intervals resample traffic phases after averaging session/quantum
repetitions within each phase. They are conditional on the fixed
calibration from two separate phases.

\begin{figure*}[t]
\centering
\includegraphics[width=\textwidth]{\quclevalpath/figures/fig5_6_execution_abstraction.pdf}
\caption{Model errors relative to Full Coupling in the separate swapping
workload, at a 0.8-ms interval. Models are adjacent within each load.
(a) Latency MAE. (b) Expected-fidelity bias in percentage points,
$100(\mathbb{E}[F_{\rm model}]-\mathbb{E}[F_{\rm Full}])$, on one common
zero-based axis. Decoupled Timing is absent from (b) because it has no
state prediction. Whiskers are 95\% traffic-cluster bootstrap intervals.}
\label{fig:abstraction}
\end{figure*}

At load 0.7, No R Contention produces \quclNoQueueMAEHigh~ms latency MAE
and expected-fidelity bias $+\quclNoQueueBiasHigh$ (about +6.06 percentage
points). Fixed Classical Delay and Decoupled Timing have smaller but
nonzero MAEs of \quclFixedMAEHigh and \quclDecoupledMAEHigh~ms.
Fixed Classical Delay's fidelity bias changes from
$\quclFixedBiasMid$ at load 0.35 to $\quclFixedBiasHigh$ at load 0.7:
packet simplification is not uniformly optimistic. As a control, at the
2-ms interval R waiting disappears and No R Contention has zero observed
latency and fidelity error. Thus, the error depends on whether the removed
execution dependency is active.
